\chapter{Открытый банкинг, A2A и ISO 20022}\label{ch:open-banking}
Когда покупатель подтверждает платёж в~приложении своего банка, а~не
вводит PAN, инженерная модель меняется целиком: в~центре оказываются
счёт, право доступа к~нему и договорённости на уровне API. Глава
разбирает, как устроен открытый банкинг на уровне регулирования,
платёжной системы и~языка сообщений.

\section{PSD2 и SCA: зачем банку открывать доступ к~счёту}\label{sec:psd2-sca}

PSD2 перевела доступ к~счёту из частного банковского канала
в~регулируемый интерфейс. Директива описывает не только платёж как
денежное событие, но и~роли, через которые этот платёж теперь
разрешается инициировать. В~её словаре есть \textbf{PSU (payment
service user)} -- клиент, дающий согласие, \textbf{ASPSP (account
servicing payment service provider)} -- банк или другой провайдер,
ведущий счёт, и лицензируемые третьи стороны: \textbf{AISP} для
доступа к~данным счёта и \textbf{PISP} для инициирования
платежа~\cite{psd2}. Банк счёта больше не может считать свой
интернет-банк единственной допустимой точкой входа в~операцию.

Если платёж инициирован через PISP, директива требует от него сразу
после инициации дать плательщику подтверждение и ссылочный
идентификатор операции. Банку счёта PISP должен передать ссылочный
идентификатор этой же операции~\cite{psd2}. Для интегратора это значит: объект
\enquote{инициация платежа} с~самого начала имеет собственный
идентификатор, а~не растворяется в~пользовательском интерфейсе.

SCA накладывается на этот режим как отдельный уровень контроля. Банк
счёта убеждается, что плательщик подтверждает доступ к~счёту
или платёж как минимум двумя независимыми элементами из трёх
категорий: знание, владение, биометрия~\cite{eba-rts-sca}. В~карточном
мире похожую задачу частично решают 3-D Secure и собственные
антифрод-механизмы эмитента. В~счётном мире SCA -- базовое условие
самого доступа к~счёту.

Для удалённого электронного платежа этого недостаточно. RTS требуют
dynamic linking: до подтверждения плательщик видит сумму
и получателя, код аутентификации привязан именно к~этой сумме
и~этому получателю, любое изменение одного из двух параметров
аннулирует ранее сгенерированный код~\cite{eba-rts-sca}.
Подтверждается не абстрактный \enquote{вход в~приложение},
а конкретная операция с~конкретными денежными атрибутами.

SCA связывает между собой согласие клиента, банк счёта и~сам платёж.
Когда продукт показывает пользователю один счёт, а~в~банковском
приложении для подтверждения всплывает другой получатель или другая
сумма, это не косметический дефект интерфейса, а~разрыв модели
доверия.

RTS одновременно перечисляют ситуации, в~которых SCA можно
не применять каждый раз. В~список входят малая сумма, доверенный
получатель, повторяющаяся операция на одну и~ту же сумму одному
и~тому же получателю, перевод между собственными счетами у~одного
провайдера и анализ транзакционного риска (TRA). Рядом живут
операции, инициированные продавцом, но они не сводятся к~тому же
набору исключений~\cite{eba-rts-sca}. Для инженера это разделение
двух классов поведения: обязательного подтверждения и запрошенного
упрощения.

\importantnote{%
  Исключение нельзя навязать банку извне: третья сторона может только
  запросить упрощённый сценарий, а~решение всегда остаётся за ASPSP,
  который ведёт счёт плательщика.
}

PSD2 и RTS задают жёсткую схему: банк счёта остаётся владельцем
аутентификации, согласия и окончательного решения по операции.
В~карточной системе продавец видит только эквайера и PSP. В~открытом
банкинге ASPSP становится явным участником продукта. Он управляет
экраном согласия, проводит SCA и возвращает не только допуск,
но и~отказ как часть нормального протокольного поведения.

Банки сопротивлялись PSD2 не из консерватизма, а~по экономической
причине: директива каннибализирует их собственный продукт. До PSD2
банк контролировал и счёт, и интерфейс доступа к~нему. Каждый платёж
через интернет-банк или мобильное приложение генерировал данные,
которые банк мог монетизировать (cross-sell, кредитный скоринг,
аналитика). PISP забирает инициацию платежа в~свой интерфейс, AISP --
доступ к~данным счёта. Банк продолжает вести счёт и~нести расчётный
риск, но теряет контакт с~клиентом в~момент платежа и~уступает
агрегатору данные, которые раньше были эксклюзивными. Отвергнутая
альтернатива -- добровольный API без регуляторного принуждения --
не работала: до PSD2 банки годами игнорировали запросы финтехов
на доступ, и Еврокомиссия сочла принуждение единственным способом
создать конкурентный рынок~\cite{psd2}.

\section{Открытый банкинг: AISP, PISP и~поток согласия}\label{sec:open-banking}

Когда счёт становится доступным через регулируемый интерфейс,
\textbf{AISP} и \textbf{PISP} расходятся на уровне доменной модели.
AISP работает с~доступом к~данным счёта и~истории операций, тогда как
PISP отвечает за инициацию конкретного платежа~\cite{psd2}. В~карточном
мире обе роли часто спрятаны внутри одного PSP, а~в~открытом банкинге
их приходится хранить и проектировать раздельно.

Продукт не может писать новую интеграцию под каждый банк, поэтому
поверх регуляторной рамки появились API-профили. В~Великобритании
эту роль выполняют стандарты Open Banking UK, в~континентальной
Европе -- Berlin Group NextGenPSD2~\cite{openbanking-uk,berlin-group}.
Они описывают один и~тот же класс задач на немного разных уровнях.
Open Banking UK публикует Read/Write API-профили, security profile
и customer experience guidelines. Berlin Group описывает общий XS2A
framework с~моделями consent, вариантами SCA и~набором платёжных
продуктов. Для интегратора это не академическое различие: оно
определяет, где искать истину о~поведении банка -- в~профиле API,
в security profile или в~документе о~пользовательском потоке.

\begin{figure}[tbp]
  \centering
  \includefiguresvg[width=.45\textwidth]{assets/figures/ch19-open-banking/pisp-flow}
  \caption{Типовой поток PISP в~открытом банкинге.}\label{fig:pisp-flow}
\end{figure}

В~реальной интеграции согласие и~платёж почти никогда не один
и~тот же объект. На стороне AISP типовой поток начинается с~создания
ресурса consent. В Open Banking UK AISP создаёт
\texttt{account-access-consent} с~использованием \textit{client
credentials grant}, получает \texttt{ConsentId} и начальный статус
\texttt{AwaitingAuthorisation}. После аутентификации пользователя
у~банка статус меняется на \texttt{Authorised} или \texttt{Rejected};
в~профиле версии v3.1.x отзыв согласия исторически живёт как
отдельное состояние \texttt{Revoked}, но к~ресурсам, созданным
начиная с v3.1.4, статус \texttt{Revoked} не применяется --
их корректное конечное состояние --
\texttt{Rejected}~\cite{openbanking-uk-account-access-consents}.
Доступ к~данным счёта -- собственный жизненный цикл, а~не побочный
эффект входа пользователя в~приложение.

У PISP тот же принцип проявляется ещё сильнее. В~типовом OB UK-потоке
PISP сначала создаёт ресурс \texttt{payment-consent} и~получает
\texttt{ConsentId}, после чего пользователь уходит в ASPSP для
аутентификации и подтверждения, и~только после этого PISP создаёт
собственно платёжный ресурс. Для consent в Open Banking UK заданы
статусы \texttt{AwaitingAuthorisation}, \texttt{Rejected},
\texttt{Authorised} и \texttt{Consumed}; \texttt{Consumed} означает
лишь, что согласованное действие было использовано, а~не то, что
деньги уже окончательно дошли до
получателя~\cite{openbanking-uk-payment-initiation-profile,openbanking-uk-domestic-payment-consents}.
Команда, привыкшая к~экранному \enquote{платёж подтверждён}, часто
путает эти статусы.

\begin{figure}[tbp]
  \centering
  \includefiguresvg[width=.75\linewidth]{assets/figures/ch19-open-banking/consent-lifecycle}
  \caption{Жизненный цикл consent в Open Banking UK: статусы PISP и AISP.}\label{fig:consent-lifecycle}
\end{figure}

На уровне аутентификации Open Banking UK поддерживает и redirect,
и развязанный подход. Причина техническая: redirect жёстко привязан
к~устройству и каналу, плохо переносится на сценарии вроде телефонии,
POS или устройств без полноценного браузера. Развязанный поток
оставляет SCA у~банка, но не требует, чтобы пользователь продолжал
её в~том же клиенте, из которого пришёл
запрос~\cite{openbanking-uk-authentication-methods}. Для платёжного
продукта вопрос звучит не как \enquote{есть ли у~нас редирект},
а~как \enquote{умеем ли мы пережить смену канала и асинхронный
возврат результата}.

Отсюда набор объектов, который продукт обязан хранить явно. Минимум
включает \texttt{consent\_id}, \texttt{payment\_id}, локальный
идентификатор сессии редиректа, срок жизни токена доступа,
идентификатор операции у~банка, итоговый статус платёжного ресурса
и~момент, после которого согласие или токен уже нельзя использовать
повторно. Без разведения этих сущностей любая проблема выглядит
одинаково: \enquote{пользователь не вернулся}, хотя в~реальности это
просроченный access token, отклонённое согласие, отмена у~банка
или отложенный финальный статус.

Продуктовая разница с~карточной страницей оплаты принципиальна.
На карточной checkout-странице продавец или его PSP управляет почти
всем взаимодействием до ответа авторизации. В~открытом банкинге
пользователь проходит часть пути в~интерфейсе банка, а~продукт видит
результат только через редирект, callback или повторный запрос
статуса. Open banking плохо переносит модель \enquote{один HTTP-запрос
-- одно денежное решение} -- здесь нужен отдельный автомат состояний.

\importantnote{%
  Открытый банкинг сам по себе не движет деньги: это доступ к~счёту
  через API и управление согласием. Денежный поток начинается там, где
  средства реально переводятся между счетами.
}

\section{Открытый банкинг в~России}\label{sec:open-banking-russia}

Российская модель открытого банкинга существенно отличается
от европейской. PSD2 -- директива Европейского союза с~обязательным
характером для банков ЕС; российский подход развивается как
поэтапная инициатива Банка России и~финтех-сообщества с~акцентом
на стандартизацию открытых API и постепенный переход
от рекомендательного режима к~обязательному для определённых
сегментов рынка~\cite{cbr-open-api,afe-open-api}.

\textbf{Концепция открытых API.} 9~ноября 2022~года Банк России
опубликовал \enquote{Концепцию внедрения Открытых API на финансовом
рынке}: документ описывает целевую модель с~набором
стандартизированных API для обмена данными и инициирования операций
со счетами, едиными правилами информационной безопасности
и поэтапным расширением круга участников
и сценариев~\cite{cbr-open-api-concept}.

\textbf{Стандартизация и роль АФТ.} Технические стандарты открытых
API в~России разрабатываются при участии Ассоциации развития
финансовых технологий (АФТ) \enquote{ФинТех} -- профессионального
сообщества, объединяющего банки, финтех-компании и~Банк России.
АФТ публикует рекомендованные стандарты OpenAPI для разных сценариев:
получение баланса и~выписки, инициирование платежа, аутентификация
клиента~\cite{afe-open-api}.\footnote{Состав опубликованных АФТ
стандартов и их статус (рекомендованные, пилотные, обязательные)
сверяются по сайту АФТ и~Банка России перед публикацией главы.}

\textbf{Связь с~СБП и~СБПэй.} В~отличие от европейской модели, где
PSD2 и SEPA Instant -- две отдельные инициативы с~разной логикой
развития, российский открытый банкинг изначально проектируется
в~связке с~СБП. Платёжная инициация через открытый API
в~значительной части сценариев означает запуск перевода через СБП
с~уже выработанным согласием клиента; сервис СБПэй
(раздел~\ref{sec:sbp-sbpay}) -- ранний пример продуктовой реализации
этой модели для розничных сценариев.

\textbf{Различие моделей.} PSD2 в~ЕС -- регуляторное обязательство
банков предоставить доступ к~счёту по запросу TPP (с~авторизацией
клиента). Российская модель идёт от пилотов и рекомендаций
к~обязательным API поэтапно: сначала добровольное участие банков,
потом обязательное -- для определённых типов API и определённых
участников. Право клиента отказать в~доступе сохраняется в~обоих
подходах, но архитектура внедрения различается по уровню
регуляторного давления и~по темпу.

\importantnote{%
  Применять европейскую терминологию AISP и PISP к~российской
  модели нужно с~оговоркой: это полезные функциональные роли,
  но в~российской регуляторике они не закреплены как
  отдельные лицензируемые статусы. Право инициировать платёж
  и~читать данные по счёту в~России опирается на договорные
  отношения участника с~банком и~на стандартизированный
  механизм согласия, а~не на отдельную лицензию TPP по
  PSD2-модели.

}

\section{ФАПИ.СЕК: российский профиль безопасности открытых API}\label{sec:fapi-sec}

Концепция Банка России задаёт регуляторную рамку, но не отвечает
на вопрос, по какому протоколу две системы реально согласуют
доступ к~счёту. Этот вопрос закрывают два технических уровня
поверх неё. Первый~-- международный профиль безопасности FAPI
(Financial-grade API), на котором построены открытый банкинг
Великобритании и~часть европейских внедрений. Второй~--
российская адаптация FAPI к~рамке информационной безопасности
Банка России и~отечественной криптографии, оформленная как
СТО БР ФАПИ.СЕК~\cite{fapi-sec-2024}.

\subsection*{FAPI как профиль безопасности}

Базовый OAuth~2.0 проектировался как фреймворк авторизации
доступа к~произвольным API, и~его настроек самих по~себе
недостаточно для финансового сценария. Параметры запроса
едут в~URL, ответ авторизационного эндпоинта возвращается
без подписи, токен доступа не~привязан ни~к~ключу клиента,
ни~к~TLS-сессии. На уровне фреймворка одна и~та же связка работает
и~для входа на новостной сайт через Google, и~для авторизации
платежа со~счёта клиента. Этой симметрии финансовый рынок
позволить себе не~может.

Профиль FAPI разрабатывает рабочая группа OpenID Foundation. Он
поднимает базовый OAuth~2.0 и~OpenID Connect до уровня, пригодного
для операций с~деньгами. \textbf{FAPI 1.0 Part~2: Advanced}
(финализирован 12~марта 2021~года) обязал использовать подписанный
JWT request object вместо параметров в~URL, sender-constrained
токены через mTLS, JARM для подписанного ответа авторизационного
эндпоинта, аутентификацию клиента через \texttt{tls\_client\_auth}
или \texttt{private\_key\_jwt} и~PKCE с~алгоритмом S256. Публичные
клиенты в~этом профиле не~поддерживаются~\cite{fapi-1-part2}.

\textbf{FAPI~2.0 Security Profile} (финализирован 22~февраля
2025~года, 82~голоса за, ноль против) сделал ещё несколько шагов.
Pushed Authorization Requests (PAR, RFC~9126~\cite{rfc-9126-par})
стал обязательным: авторизационный запрос больше не~передаётся
через URL даже в~подписанном виде, а~напрямую отправляется
на отдельный эндпоинт авторизационного сервера, который возвращает
короткий \texttt{request\_uri}. Sender-constrained токены остались
обязательными, но к~mTLS добавилась альтернатива в~виде DPoP
(RFC~9449~\cite{rfc-9449-dpop}) для случаев, где взаимный TLS
неудобен инфраструктурно. PKCE S256 и~TLS~1.2+ остались
обязательными~\cite{fapi-2-security-profile}.

\textbf{FAPI~2.0 Message Signing} (финализирован 26~сентября
2025~года) дополнил профиль обязательной подписью самих сообщений
запроса и~ответа на ресурсном уровне, а~не~только на этапе
авторизации. Это нужно для non-repudiation в~сценариях, где
финансовый институт обязан доказать содержание операции спустя
длительное время после её исполнения~\cite{fapi-2-message-signing}.

Профили FAPI описывают \emph{что} должны делать стороны и~\emph{как}
выглядит сообщение на проводе. \emph{Чем} оно подписывается~--
RSA, ECDSA, ГОСТ,~-- профиль сам по~себе не~предписывает; он
ссылается на алгоритмы JWS из~IANA-реестра. Поэтому национальная
адаптация FAPI~-- это не~переписывание логики протокола, а~подмена
криптографической части и~обвязка из~требований местного регулятора.

\subsection*{ФАПИ.СЕК: что наследуется, что меняется}

СТО~БР ФАПИ.СЕК-1.6-2024, принятый и~введённый в~действие приказом
Банка России от~7~октября 2024~года №~ОД-1615, заменил редакцию
2020~года и~стал актуальным профилем безопасности для российских
открытых API. Стандарт разработан Ассоциацией ФинТех совместно
с~ИнфоТеКС при участии Банка
России~\cite{fapi-sec-2024,cbr-fapi-sec-9908}.

Структура стандарта прямо наследует профили FAPI. В~ФАПИ.СЕК
выделены три раздела: профиль безопасности OpenID API в~режиме
только для чтения (соответствует FAPI 1.0~Part~1: Baseline);
профиль безопасности в~режиме чтения и~записи (соответствует
FAPI 1.0~Part~2: Advanced) и~JARM как защищённый JWT-режим ответа
авторизационного эндпоинта~\cite{fapi-sec-2024,fapi-1-part1,fapi-1-part2,oidf-jarm}.
Содержательно расширенный профиль ФАПИ.СЕК требует тех же
механизмов, что и~FAPI Advanced: подписанный request object,
взаимный TLS на транспортном уровне, sender-constrained токены,
\texttt{private\_key\_jwt} либо \texttt{tls\_client\_auth}, PKCE~S256.

Точка сдвига проходит по~криптографии. ФАПИ.СЕК предписывает
исключительное применение отечественных криптографических
алгоритмов согласно методическим рекомендациям ТК~26
МР~26.2.002-2024 «Использование российских криптографических
алгоритмов в~протоколах OpenID Connect», утверждённым
19~ноября 2024~года~\cite{tc26-mr-26-2-002-2024}. Подпись
JWT и~проверка подписи~-- по~ГОСТ Р~34.10-2012 на эллиптических
кривых вместо RSA или ECDSA. Хеш-функция в~JWS, JWE и~JWT~--
ГОСТ Р~34.11-2012 «Стрибог» вместо SHA-2. Симметричное шифрование
в~JWE~-- по~ГОСТ Р~34.12-2015 («Кузнечик» или «Магма») вместо
AES. Транспорт~-- TLS с~российскими криптонаборами (ГОСТ TLS),
а~не~обычный TLS~1.2/1.3 с~международными шифросьютами. Подпись
запросов и~хранение ключей выполняются исключительно средствами
сертифицированных СКЗИ; мягкие реализации криптографии в~ФАПИ.СЕК
не~допускаются.

\begin{figure}[tbp]
  \centering
  \includefiguresvg[width=\linewidth]{assets/figures/ch19-open-banking/fapi-sec-stack}
  \caption{Стек ФАПИ.СЕК: процедурный слой FAPI 1.0~Advanced
  и~JARM наследуется без~изменений; международная криптография
  и~международные TLS-криптонаборы заменяются на отечественные
  по~МР~26.2.002-2024 ТК~26.}\label{fig:fapi-sec-stack}
\end{figure}

Отдельный стандарт СТО~БР ФАПИ.ПАОК-1.0-2024, принятый тем же
приказом, описывает профиль аутентификации, инициированной
клиентом по~отдельному каналу (Client Initiated Backchannel
Authentication, CIBA)~\cite{fapi-paok-2024}. Сценарий нужен там,
где у~устройства, инициирующего доступ, нет полноценного
браузера или экрана для подтверждения: оплата с~умной колонки,
подтверждение через мобильное банковское приложение в~ответ
на инициацию из~контактного центра, аутентификация на устройстве,
не~имеющем входа клиента. ПАОК относится к~ФАПИ.СЕК как decoupled
authentication~-- к~основному redirect-сценарию у~Open Banking
UK~\cite{openbanking-uk-authentication-methods}: разные каналы
доставки шага аутентификации, общая модель доверия.

\subsection*{Согласие, аутентификация и подпись запроса}

ФАПИ.СЕК наследует от FAPI разделение трёх сущностей, которые
в~продукте легко склеиваются в~одну. Первое~-- объект согласия,
живущий на стороне ASPSP с~собственными состояниями (создано,
ожидает аутентификации клиента, активно, отозвано, истекло).
Второе~-- сама аутентификация клиента, выполняемая средствами
банка по~уровню, согласованному с~регуляторными требованиями.
Третье~-- подписанный запрос инициации операции, в~котором
клиентское приложение доказывает, что согласие действительно
и~операция совпадает с~тем, что подтвердил клиент.

Типовой сценарий инициации платежа выглядит так. Клиентское
приложение формирует request object~-- JWT с~описанием операции,
подписанный закрытым ключом клиента средствами СКЗИ по~ГОСТ
Р~34.10-2012. Object отправляется на PAR-эндпоинт банка и~получает
в~ответ короткий \texttt{request\_uri}. Браузер клиента
перенаправляется в~банк по~этому \texttt{request\_uri}; банк
проверяет подпись запроса, выполняет аутентификацию клиента
(multi-factor по~требованиям к~средствам идентификации), показывает
экран согласия с~параметрами операции и~возвращает
\texttt{authorization\_code} по~callback. Клиентское приложение
меняет \texttt{authorization\_code} на \texttt{access\_token},
привязанный к~mTLS-сертификату или DPoP-ключу, и~использует токен
для самого вызова API инициации платежа.

\begin{figure}[tbp]
  \centering
  \includefiguresvg[width=\linewidth]{assets/figures/ch19-open-banking/fapi-sec-flow}
  \caption{Поток инициации платежа через ФАПИ.СЕК с~расчётом
  через ОПКЦ СБП. Шаги~2--3 и~7~следуют процедуре FAPI~2.0
  Security Profile (PAR, sender-constrained tokens), подпись
  и~хранение ключей выполняются средствами сертифицированных
  СКЗИ.}\label{fig:fapi-sec-flow}
\end{figure}

Каждый шаг этой цепочки~-- отдельная точка отказа со~своим
кодом ответа. Подпись request object не~сошлась~-- ФАПИ.СЕК
требует отклонить запрос до~показа экрана согласия. Согласие
истекло или отозвано клиентом~-- \texttt{access\_token}
отзывается, банк отказывает в~инициации со~статусом, отличным
от~\enquote{недостаточно средств} или \enquote{лимит превышен}.
Сертификат клиента в~mTLS не~совпал с~привязкой токена~--
повторно ходить за тем же токеном бессмысленно, нужно начинать
с~PAR заново. Продукт, который сводит все эти отказы в~одну
колонку \enquote{платёж не~прошёл}, теряет управляемость точно
так же, как описано в~\ref{sec:open-banking-boundaries}
для Open Banking UK.

\begin{figure}[tbp]
  \centering
  \includefiguresvg[width=.85\linewidth]{assets/figures/ch19-open-banking/fapi-sec-consent-lifecycle}
  \caption{Жизненный цикл объекта согласия по~разделу~9.3.3
  ConsentStatus стандарта СТО~БР \enquote{Открытые программные
  интерфейсы. Получение согласия на~доступ к~банковским счетам
  Пользователя} (приказ Банка России от~19.12.2025
  № ОД-2895). DELETE-операция доступна как до~авторизации
  пользователем, так и~после~неё.}\label{fig:fapi-sec-consent-lifecycle}
\end{figure}

Отношение ФАПИ.СЕК к~стандартам аутентификации в~карточном
платеже отдельное. SCA в~3-D~Secure (см.~\ref{ch:3ds})
аутентифицирует держателя карты по~конкретной операции и~работает
внутри карточной сети. ФАПИ.СЕК аутентифицирует клиента банка
по~доступу к~счёту и~работает поверх инфраструктуры открытых
API. У~двух механизмов общий принцип сильной аутентификации
из~двух факторов разных категорий, но~разные точки доверия
и~разные операционные потоки. Для пользователя они могут выглядеть
одинаково (push-уведомление в~приложении банка), но~регулируются
и~сертифицируются раздельно.

\subsection*{Связка с~СБП и~российский график обязательности}

ФАПИ.СЕК~-- не~расчётная инфраструктура, а~технический протокол
доступа к~счёту и~инициации операций над ним. Когда инициация
платежа доходит до этапа исполнения, расчёт идёт через
ту же инфраструктуру, что и~любой другой A2A-платёж в~России~--
через ОПКЦ СБП и~платёжную систему Банка
России~(\ref{sec:sbp-protocol}). Это и~есть архитектурное место
ФАПИ.СЕК: открытый банкинг через регулируемый интерфейс
к~счёту, расчёт через мгновенную межбанковскую систему. Ни
ФАПИ.СЕК без СБП, ни~СБП без стандартизованного API доступа
к~счёту не~дают полноценной счётной модели платежа.

Концепция Банка России от~9~ноября 2022~года задала поэтапный
график движения от~рекомендательного режима к~обязательному
для отдельных API в~четырёх секторах: банковском, инвестиционном,
страховом, микрофинансовом~\cite{cbr-open-api-concept}. По~банковскому
сектору рекомендательное использование API \enquote{Информация
об~организации} объявлено с~2023~года, обязательное~-- с~2024~года;
для API сведений о~счетах и~продуктах клиента, а~также для
инициирования платежа~-- рекомендательное с~2024~года, обязательное
с~2025~года. Страховой сектор движется с~годовой задержкой
относительно банковского.

24~декабря 2025~года Банк России опубликовал обновлённый комплекс
из~десяти стандартов Открытых API, утверждённых приказами Банка
России от~19~декабря 2025~года. Стандарты сохраняют рекомендательный
характер; переход участников пилотных проектов на новые версии
спецификаций по~счетам физических и~юридических лиц запланирован
на 1~октября 2026~года~\cite{cbr-open-api-update-2025}. Окончательные
сроки и~состав обязательных API меняются по~мере доработки концепции;
точные даты сверяются по~сайту Банка России и~Ассоциации
ФинТех~\cite{cbr-open-api,afe-open-api}.

\subsection*{Российский профиль рядом с~европейскими}

Поставив ФАПИ.СЕК рядом с Open Banking UK, Berlin Group NextGenPSD2
и базовым FAPI 2.0, видно, по каким осям российский профиль
действительно отличается от европейских.

\begin{table}[tbp]
  \begingroup
  \scriptsize
  \setlength{\tabcolsep}{4pt}
  \renewcommand{\arraystretch}{1.2}
  \makebox[\textwidth][l]{%
    \begin{tabularx}{\fullwidth}{
        @{}>{\RaggedRight\arraybackslash\bfseries}p{0.18\fullwidth}
        >{\RaggedRight\arraybackslash}X
        >{\RaggedRight\arraybackslash}X
        >{\RaggedRight\arraybackslash}X
        >{\RaggedRight\arraybackslash}X@{}
      }
      \toprule
      \rowcolor{Panel}
      Параметр & ФАПИ.СЕК (РФ) & Open Banking UK & Berlin Group (ЕС) & FAPI 2.0 (OIDF) \\
      \midrule
      \rowcolor{Soft}
      Базовый профиль & FAPI 1.0 Advanced + ГОСТ-надстройка & FAPI 1.0 Advanced & собственный профиль на FAPI и OAuth 2.0 & FAPI 2.0 \\
      Криптография & ГОСТ Р 34.10/34.11/34.12 в~JWS, JWE, JWK & RSA, ECDSA, SHA-2 & RSA, ECDSA, SHA-2 & задаётся реализующим профилем \\
      \rowcolor{Soft}
      Транспорт & ГОСТ TLS & TLS 1.2+ & TLS 1.2+ & TLS 1.2+ \\
      Подпись запроса & request object на ГОСТ, СКЗИ обязательно & signed JWT, mTLS & signed JWT, mTLS либо QSeal & PAR обязательно \\
      \rowcolor{Soft}
      Лицензируемые роли & нет; договор с~банком + соответствие стандарту & AISP, PISP по PSD2 & AISP, PISP, CBPII по PSD2 & не определяет \\
      Расчётный канал & ОПКЦ СБП & Faster Payments & SEPA Instant либо локальная RTGS & не определяет \\
      \rowcolor{Soft}
      Статус & рекомендательный, поэтапный переход к~обязательному & обязательный для девяти крупнейших банков & обязательный по PSD2 & добровольная адаптация \\
      \bottomrule
    \end{tabularx}%
  }
  \endgroup
  \caption{ФАПИ.СЕК на одном полотне с~европейскими профилями
  открытого банкинга.}\label{tab:fapi-sec-comparison}
\end{table}

Главный сдвиг ФАПИ.СЕК относительно Open Banking UK и~Berlin Group~--
не~процедурный (его процедуры почти полностью наследуются от~FAPI),
а~инфраструктурный: подмена международной криптографии и~TLS
на отечественные алгоритмы, требование сертифицированных СКЗИ
у~всех участников, отсутствие лицензируемых TPP-ролей в~пользу
договорной модели с~банком. Для команды, проектирующей продукт
на ФАПИ.СЕК, FAPI 1.0~Advanced и~JARM остаются основным источником
понимания процедурной части, а~МР~26.2.002-2024 ТК~26 и~821-П
с~851-П (см.~\ref{sec:regulation-cbr-provisions})~-- источниками
криптографических и~операционных требований.

Состав обязательных API, тарифы пилотов, точные сроки перехода
и~актуальная редакция приказов меняются и~сверяются по~сайту Банка
России и~публикациям Ассоциации ФинТех; в~книге зафиксирована
архитектурная рамка, а~не~оперативная справка.

\section{A2A-платежи как отдельная система}\label{sec:a2a-payments}

Когда PISP инициирует перевод со счёта клиента, встаёт вопрос,
по какой расчётной системе реально пойдут деньги. Здесь появляется
\textbf{A2A (account-to-account)}: если open banking даёт способ
получить согласие и безопасно обратиться к~счёту, то A2A задаёт саму
модель движения денег между банковскими счетами. Эти слои часто идут
вместе, но это не одно и~то же.

Разные юрисдикции собирают A2A по-разному. Pix в~Бразилии стал
розничной системой мгновенных переводов, работающей с 16 ноября
2020 года~\cite{pix}. По первоначальной редакции правил (Resolução
BCB № 1/2020) обязательное участие распространялось на институты
с~более чем 500~тыс.\ активных счетов; в~ноябре 2024 года эти правила
были скорректированы Резолюцией BCB № 429/2024, актуальный порог
сверяется по действующей редакции~\cite{pix-participants}. У UPI
в~Индии другая конструкция: единая инфраструктура NPCI сочетается
с~конкуренцией банковских и небанковских приложений поверх
неё~\cite{upi}. В~Европе эталоном остаются мгновенные кредитовые
переводы SEPA: деньги доступны получателю в~течение секунд, 24 часа
в~сутки и 365 дней в~году~\cite{sepa-instant}. СБП из предыдущей
главы -- локальная российская мгновенная расчётная система.

SPI -- центральная инфраструктура расчёта мгновенных платежей
Banco Central do Brasil под Pix: валовой расчёт в~реальном времени
через специальные счета участников в~центральном банке, после расчёта
перевод становится безотзывным~\cite{pix-spi}. Pix -- не ускоренная
карточная система, а~иной тип денежного события: положительный
статус близок к~самому денежному переводу, а~не к~разрешению, как
в~карточной авторизации. Сравнение карточной и A2A-моделей по
финальности, возвратам и~спорам подробно разобрано
в~\ref{sec:sbp-vs-cards}.

\begin{table}[tbp]
  \begingroup
  \small
  \setlength{\tabcolsep}{4pt}
  \renewcommand{\arraystretch}{1.2}
  \makebox[\textwidth][l]{%
    \begin{tabularx}{\fullwidth}{
        @{}>{\RaggedRight\arraybackslash\bfseries}p{0.18\textwidth}
        >{\RaggedRight\arraybackslash}X
        >{\RaggedRight\arraybackslash}X
        >{\RaggedRight\arraybackslash}X@{}
      }
      \toprule
      \rowcolor{Panel}
      Система & Что даёт первый положительный ответ &
      Когда деньги считаются окончательными &
      Как устроены возврат и~спор \\
      \midrule
      \rowcolor{Soft}
      Карточный платёж &
      Разрешение на авторизацию и резервирование средств &
      После отдельного клиринга и~расчёта по правилам схемы &
      Возврат живёт отдельно от авторизации/подтверждения; возвратный платёж встроен
      в~сетевые правила \\
      A2A-платёж &
      Обычно подтверждение инициирования перевода или его исполнения,
      в~зависимости от системы &
      Когда межбанковский перевод исполнен и~система считает его окончательным &
      Возврат -- новый push-перевод; универсального возвратного платежа
      уровня карточной схемы обычно нет \\
      \bottomrule
    \end{tabularx}%
  }
  \endgroup
  \caption{Карточная и A2A-системы расходятся по lifecycle, а~не только
  по пользовательскому интерфейсу.}\label{tab:card-vs-a2a}
\end{table}

Особенно заметна разница в~электронной коммерции. Для продавца
мгновенная расчётная система убирает длинный разрыв между
\enquote{покупатель подтвердил} и \enquote{деньги дошли}, и~вместе
с~этим разрывом исчезает запас на ошибку. Если заказ отгружается
на основании промежуточного статуса или возврат оформляется так, будто
у~системы есть штатный возвратный платёж, внутренний учёт быстро
расходится с~реальным движением денег. A2A требует другой модели
управления заказом, другой поддержки и другой логики сверки.

\section{ISO 20022 вне карточного мира}\label{sec:iso20022}

ISO 20022 задаёт семантическую модель: у~поля есть имя и~место
в~общей бизнес-структуре сообщения. ISO 20022 удобен там, где один
и~тот же платёж должен одновременно быть понятен маршрутизации,
санкционному и AML-контролю, сверке, выписке
и отчётности~\cite{iso-20022}.

Три семейства сообщений соответствуют разным участкам цепочки.
\texttt{pain} покрывает \textit{customer initiation} -- запросы
клиента на инициирование платежа; \texttt{pacs} используется
в \textit{payments clearing and settlement} между финансовыми
институтами; \texttt{camt} относится к \textit{cash management}
и~отвечает за статусы, выписки, уведомления и связанные сообщения
управления ликвидностью~\cite{iso-20022-business-areas}. Эти префиксы
описывают разные точки зрения на одну и~ту же операцию.

В~практической цепочке клиент или его система
формирует \texttt{pain.001} как инструкцию на кредитовый перевод.
На межбанковском уровне перевод идёт как \texttt{pacs.008}; ответ
о~состоянии на этом же уровне часто приходит через \texttt{pacs.002}.
Когда деньги отражаются в~учёте клиента, банк присылает
\texttt{camt.054} или включает операцию в~выписку вида
\texttt{camt.053}. У~каждого шага своя аудитория и свой момент
истины. Если приложение видит только пользовательский API, а~сверка
живёт по \texttt{camt}, между ними нужно явное внутреннее
соответствие.

Богатая структура сообщения окупается на сверке. Для сверки
и расследования нужны сумма, валюта и~поля вроде \texttt{EndToEndId}:
идентификаторы дебитора и кредитора, структурированное назначение
платежа и~код причины в~статусном сообщении. В~карточном мире часть
этих данных живёт во внешнем процессинговом диалекте (\textit{processing}) или в~файлах
клиринга. В ISO 20022 они закладываются в~саму модель сообщения.

Стоимость миграции связана не только с XML-схемой. При переходе
на новый профиль ISO 20022 меняются правила заполнения полей, usage
guidelines, допустимые статусы и обязательность отдельных атрибутов.
Без канонической модели платежа внутри продукта такой переход
превращается в~переписывание API, оркестрации и reconciliation сразу.

Трансграничные расчёты это подтверждают наглядно. В SWIFT
coexistence-period MT и ISO 20022 для межбанковских платёжных
инструкций завершился в~ноябре 2025 года (cutover weekend
21--23 ноября)~\cite{swift-iso20022-milestone,swift-iso20022-implementation}.
После этой даты нельзя рассчитывать, что старый FIN MT сам доедет
до адресата как основной рабочий формат. С 1 января 2026 года
contingency processing для MT-отправителей и in-flow translation
тарифицируются отдельно. MT~199 и MT~299 для exception handling
сохраняются после ноября 2025 года как переходный канал и~подлежат
окончательной отмене в~ноябре 2027 года, когда их функцию полностью
переймут сообщения camt.110 и camt.111~\cite{swift-iso20022-implementation}.
Для инженера отсюда простое правило: мало знать название сообщения.
Нужно знать профиль, правила переходного периода и~ветку
исключительной обработки.

В~российской системе движение то же. Банк России ведёт отдельный
раздел по стандарту ISO 20022, описывает его как основу собственных
стандартов финансовых сообщений и отдельно поддерживает обмен
сообщениями ISO 20022 в~СПФС~\cite{cbr-iso20022-album,cbr-spfs-iso20022}.
Разговор о~богатых структурированных данных относится не только
к SWIFT и cross-border -- он касается и внутренних платёжных систем,
если они строятся вокруг счёта, а~не вокруг PAN.

\section{Чем открытый банкинг отличается от A2A и
карточной страницы оплаты}\label{sec:open-banking-boundaries}

В~разговорах о~счётных платежах постоянно склеиваются три уровня.
Первый -- доступ к~счёту и управление согласием. Второй -- расчётный
канал, по которому реально идут деньги. Третий -- язык межсистемного
сообщения. Пока эти три уровня не разведены, любой разговор
об \enquote{оплате со счёта} остаётся грубым.

\begin{table}[tbp]
  \begingroup
  \small
  \setlength{\tabcolsep}{4pt}
  \renewcommand{\arraystretch}{1.2}
  \makebox[\textwidth][l]{%
    \begin{tabularx}{\fullwidth}{
        @{}>{\RaggedRight\arraybackslash\bfseries}p{0.18\textwidth}
        >{\RaggedRight\arraybackslash}X
        >{\RaggedRight\arraybackslash}X
        >{\RaggedRight\arraybackslash}X@{}
      }
      \toprule
      \rowcolor{Panel}
      Уровень & На какой вопрос отвечает & Кто обычно им управляет &
      Что ломается, если смешать его с~соседним \\
      \midrule
      \rowcolor{Soft}
      Access / consent API &
      Кто и~на каких правах может обратиться к~счёту &
      ASPSP, API-профиль, security profile, consent management &
      Команда принимает истечение токена или отказ согласия за сбой
      расчётной системы \\
      Расчётная система &
      Как деньги переводятся между счетами и~когда они окончательны &
      Платёжная система, расчётная инфраструктура, правила системы &
      Промежуточный статус принимается за расчёт; возврат проектируется
      как отмена уже исполненного перевода \\
      \rowcolor{Soft}
      Message standard &
      Как операция описывается между системами и в~отчётности &
      ISO 20022 profile, usage guidelines, локальные альбомы сообщений &
      Внешний API и внутренний реестр обязательств завязываются на один конкретный
      формат сообщения \\
      \bottomrule
    \end{tabularx}%
  }
  \endgroup
  \caption{Открытый банкинг, расчётная система и~уровень сообщений отвечают за разные части
  одной и~той же операции.}\label{tab:account-based-layers}
\end{table}

Карточная checkout-страница устроена иначе: в~её центре -- PAN,
эмитентный канал авторизации, клиринг и dispute framework карточной
схемы. В~счётной модели в~центре оказываются consent на доступ
к~счёту, инициирование банковского перевода и финальность расчётной
системы. СБП и~открытый банкинг удобно читать рядом: СБП показывает
локальную A2A-систему, эта глава -- более общий класс систем, где
доступ к~счёту и банковский перевод становятся самостоятельной
архитектурой.

Последнее разделение -- между отказом, статусом и~спором. Отказ
на этапе consent или initiation -- часть нормального протокола:
токен истёк, клиент отменил согласие, банк не принял выбранный
сценарий SCA, счёт недоступен. Статус платежа -- отдельная задача:
деньги могут быть ещё не финальны, даже если consent уже
\texttt{Consumed} или API вернул \texttt{Accepted}. Спор после
исполнения перевода -- третья задача, и~она не обязана иметь
карточный возвратный платёж как стандартную форму разрешения конфликта.
Смешав эти три режима в~одну колонку \enquote{платёж не удался},
продукт теряет управляемость.

\practicenote{%
  При проектировании счётных платежей разведите три
  внутренние сущности: объект согласия, объект денежной операции и~объект
  межсистемного сообщения. Тогда миграция на новый банк, новую расчётную систему или
  новый профиль ISO 20022 меняет только одну сущность, а~не всю систему
  сразу.
}
